\chapter{Simulación del modelo SocialScrum}
\label{anexo:simulacion-socialscrum}

Este anexo presenta una simulación que ilustra la aplicación de SocialScrum durante tres Sprints consecutivos. El escenario se basa en patrones observados en equipos reales de desarrollo, aunque los nombres y detalles específicos son ficticios. El propósito es demostrar cómo los componentes del marco ---eventos adaptados, artefactos sociales, sistema de medición y rol del Social Guide--- se integran en la práctica diaria de un equipo.

\section{Contexto del equipo}
El equipo Phoenix es un equipo Scrum de 7 personas que desarrolla una plataforma SaaS de gestión de inventarios. Tiene 18 meses de experiencia con Scrum estándar, pero ha enfrentado problemas de comunicación y colaboración que han afectado su productividad y moral, en la Tabla \ref{tab:perfil-phoenix} se presenta su perfil detallado.

\begingroup
\scriptsize
\renewcommand{\arraystretch}{1.4}
\setlength{\LTleft}{0pt}
\setlength{\LTright}{0pt}

\begin{longtable}[H]{
  >{\RaggedRight\hspace{0pt}}p{5.2cm}   % Col 1: Característica
  >{\RaggedRight\hspace{0pt}}p{10.0cm}  % Col 2: Descripción
  }

  % --- ENCABEZADO PRIMERA PÁGINA ---
  \caption{Perfil del equipo Phoenix}
  \label{tab:perfil-phoenix}                                                                        \\
  \toprule
  \textbf{Característica} & \textbf{Descripción}                                                    \\
  \midrule
  \endfirsthead

  % --- ENCABEZADO SIGUIENTES PÁGINAS ---
  \multicolumn{2}{c}{{\bfseries \tablename\ \thetable{} -- Continuación}}                           \\
  \toprule
  \textbf{Característica} & \textbf{Descripción}                                                    \\
  \midrule
  \endhead

  % --- PIE INTERMEDIO ---
  \midrule
  \multicolumn{2}{r}{\textit{Continúa en la siguiente página...}}                                   \\
  \endfoot

  % --- PIE FINAL ---
  \bottomrule
  \endlastfoot

  Nombre del equipo
                          & Equipo Phoenix                                                          \\

  Tamaño
                          & 7 personas (1 Product Owner, 1 Scrum Master, 5 desarrolladores)         \\

  Producto
                          & Plataforma SaaS de gestión de inventarios                               \\

  Modalidad
                          & Híbrida (3 días en oficina, 2 días remotos)                             \\

  Duración del Sprint
                          & 2 semanas                                                               \\

  Madurez en Scrum
                          & 18 meses practicando Scrum estándar                                     \\

  Social Guide
                          & Laura (Scrum Master con capacitación adicional; rol compartido al 50\%) \\
\end{longtable}
\endgroup

\subsection{Situación inicial}

Antes de implementar SocialScrum, el equipo presentaba los siguientes síntomas:

\begin{itemize}
  \item Velocidad inconsistente: fluctuaciones de 30\% entre Sprints sin causa técnica clara.
  \item Dos desarrolladores (Carlos y Miguel) evitaban trabajar juntos desde un conflicto no resuelto hace 4 meses.
  \item Ana, la desarrolladora más senior, concentraba el 60\% del conocimiento crítico del sistema.
  \item Las retrospectivas se habían vuelto ``ceremoniales'': respuestas genéricas, sin mejoras reales.
  \item El PO reportaba fricción frecuente al solicitar cambios de prioridad.
\end{itemize}

\subsection{Diagnóstico inicial (Sprint 0)}

Laura realizó el diagnóstico inicial usando las herramientas de SocialScrum como se muestra en la Tabla \ref{tab:health-score-inicial}.

\begingroup
\scriptsize
\renewcommand{\arraystretch}{1.4}
\setlength{\tabcolsep}{4pt}
\setlength{\LTleft}{0pt}
\setlength{\LTright}{0pt}

\begin{longtable}[H]{
  >{\RaggedRight\hspace{0pt}}p{0.18\textwidth}  % Dimensión
  >{\Centering\hspace{0pt}}p{0.15\textwidth}   % Puntuación
  >{\RaggedRight\hspace{0pt}}p{0.60\textwidth} % Observación
  }

  % --- ENCABEZADO PRIMERA PÁGINA ---
  \caption{Health Score inicial del equipo Phoenix}
  \label{tab:health-score-inicial}                                                                             \\
  \toprule
  \textbf{Dimensión} & \textbf{Puntuación}                                              & \textbf{Observación} \\
  \midrule
  \endfirsthead

  % --- ENCABEZADO SIGUIENTES PÁGINAS ---
  \multicolumn{3}{c}{{\bfseries \tablename\ \thetable{} -- Continuación}}                                      \\
  \toprule
  \textbf{Dimensión} & \textbf{Puntuación}                                              & \textbf{Observación} \\
  \midrule
  \endhead

  % --- PIE INTERMEDIO ---
  \midrule
  \multicolumn{3}{r}{\textit{Continúa en la siguiente página...}}                                              \\
  \endfoot

  % --- PIE FINAL ---
  \bottomrule
  \endlastfoot

  Apertura
                     & 4.8
                     & Los problemas se ocultan hasta el último momento                                        \\

  Compromiso
                     & 6.2
                     & Aceptable, pero es frecuente la actitud de "ese no es mi código"                        \\

  Respeto
                     & 5.1
                     & Tensión visible entre Carlos y Miguel                                                   \\

  Foco
                     & 5.5
                     & Interrupciones frecuentes y cambio constante de contexto                                \\

  Coraje
                     & 4.2
                     & Nadie menciona explícitamente el conflicto entre Carlos y Miguel                        \\

  \midrule
  \textbf{Health Score}
                     & \textbf{5.2}
                     & \textbf{Zona "en riesgo"}                                                               \\
\end{longtable}
\endgroup

Los \gls{community-smells} identificados fueron:
\begin{enumerate}
  \item \textbf{Radio Silence} (Severidad Alta): Carlos y Miguel no se comunican directamente.
  \item \textbf{Organisational Silo} (Severidad Media): Frontend y Backend trabajan aislados.
  \item \textbf{Lone Wolf} (Severidad Media): Ana trabaja sola en módulos críticos.
\end{enumerate}

\section{Sprint 1: Establecer fundamentos}

\subsection{Social Sprint Planning}
Durante el Planning, Laura facilitó la revisión de riesgos sociales:

\begin{escenario}

  \textbf{Laura (Social Guide)}:
  \begin{quote}
    "Antes de estimar, revisemos el contexto social.
    El Health Score es 5.2 y tenemos tres \gls{community-smells} activos.
    ¿Ven algún riesgo para este Sprint?"
  \end{quote}

  \textbf{Product Owner}:
  \begin{quote}
    "Necesito que Carlos y Miguel trabajen juntos en el módulo de reportes."
  \end{quote}

  \textbf{Laura (Social Guide)}:
  \begin{quote}
    "Eso representa un riesgo. Propongo que, antes de asignar esa tarea,
    realicemos una sesión de alineación de una hora.
    ¿Están de acuerdo?"
  \end{quote}

  \textbf{Resultado}:
  \begin{quote}
    El equipo acepta la propuesta y se agenda la sesión de alineación para el día 2 del Sprint.
  \end{quote}

\end{escenario}


\subsection{Social Sprint Backlog}

\textbf{Capacidad asignada}: 8 SSP de 50 puntos totales del Sprint (16\%).

\begin{itemize}
  \item Sesión de alineación entre Carlos y Miguel (3 SSP).
  \item \gls{pair-programming} rotativo entre Ana y otro desarrollador (3 SSP).
  \item Creación y adopción del canal \#social-phoenix en Slack (2 SSP).
\end{itemize}

\subsection{Ejecución del Sprint}

\paragraph{Social Daily Standup's:}

En este nuevo evento se implementó la pregunta social diaria donde los miembros del equipo respondían a una pregunta rotativa diseñada para detectar señales tempranas de fricción o malestar. En la Tabla \ref{tab:daily-sprint1} se presentan las señales detectadas durante el Sprint 1.

\begingroup
\scriptsize
\renewcommand{\arraystretch}{1.4}
\setlength{\tabcolsep}{4pt}
\setlength{\LTleft}{0pt}
\setlength{\LTright}{0pt}

\begin{longtable}{
  >{\Centering\hspace{0pt}}p{1.6cm}     % Día
  >{\RaggedRight\hspace{0pt}}p{4.6cm}   % Pregunta
  >{\RaggedRight\hspace{0pt}}p{8.9cm}   % Señal detectada
  }

  % --- ENCABEZADO PRIMERA PÁGINA ---
  \caption{Señales detectadas en Social Daily Standup – Sprint 1}
  \label{tab:daily-sprint1}                                                          \\
  \toprule
  \textbf{Día} & \textbf{Pregunta}                        & \textbf{Señal detectada} \\
  \midrule
  \endfirsthead

  % --- ENCABEZADO SIGUIENTES PÁGINAS ---
  \multicolumn{3}{c}{{\bfseries \tablename\ \thetable{} -- Continuación}}            \\
  \toprule
  \textbf{Día} & \textbf{Pregunta}                        & \textbf{Señal detectada} \\
  \midrule
  \endhead

  % --- PIE INTERMEDIO ---
  \midrule
  \multicolumn{3}{r}{\textit{Continúa en la siguiente página...}}                    \\
  \endfoot

  % --- PIE FINAL ---
  \bottomrule
  \endlastfoot

  3
               & ¿Cómo te sientes?
               & Miguel: "Tenso" (post-sesión con Carlos)                            \\

  5
               & ¿Necesitas algo?
               & Ana: "Ayuda con documentación"                                      \\

  8
               & ¿Quién te ayudó?
               & Carlos: "Miguel me explicó el módulo"                               \\
\end{longtable}
\endgroup

\paragraph{Intervención:}\mbox{}\\[-0.1\baselineskip]

A continuación se detalla la sesión de alineación entre Carlos y Miguel, facilitada por Laura, la cual fue una intervención clave para abordar el community smell en un escenario aplicado.
\begin{escenario}

  \textbf{Sesión de alineación Carlos–Miguel}

  Laura facilita una sesión de alineación de 90 minutos con el objetivo
  de reducir fricción y habilitar trabajo conjunto durante el Sprint.

  \begin{enumerate}
    \item \textbf{Contexto}. Cada participante expone su perspectiva del conflicto original,
          sin búsqueda de culpables ni juicios personales.

    \item \textbf{Impacto}. Ambos reconocen que evitar la colaboración estaba afectando
          al equipo, en particular la velocidad y la gestión de dependencias.

    \item \textbf{Acuerdos}. Se establecen tres normas de trabajo conjunto:
          \begin{itemize}
            \item Comunicación escrita para decisiones relevantes.
            \item Escalamiento inmediato a Laura ante señales de fricción.
            \item Retrospectiva privada conjunta al final del Sprint.
          \end{itemize}
  \end{enumerate}

  \textbf{Resultado}. La sesión permite habilitar la colaboración controlada
  entre Carlos y Miguel para el resto del Sprint.

\end{escenario}

\subsection{Social Sprint Retrospective}

Fase 3 (Análisis) reveló patrón: ``No decimos cuando algo nos molesta hasta que explota.''
Análisis de causa raíz (5 Por qués):
\begin{enumerate}
  \item ¿Por qué no hablamos?
        Miedo a reacciones negativas.

  \item ¿Por qué hay miedo?
        Experiencias pasadas de reacciones defensivas.

  \item ¿Por qué hubo reacciones defensivas?
        No existía un espacio seguro para dar feedback.

  \item ¿Por qué no había un espacio seguro?
        Nunca se estableció explícitamente.

  \item \textbf{Causa raíz}: Falta de normas explícitas para dar feedback crítico.
\end{enumerate}

\textbf{Compromiso}: Crear ``Acuerdo de Feedback'' con 5 reglas básicas (3 SSP para Sprint 2).

\subsection{Resultados Sprint 1}

En la Tabla \ref{tab:resultados-sprint1} se presentan los resultados cuantitativos del Sprint 1.

\begingroup
\scriptsize
\renewcommand{\arraystretch}{1.4}
\setlength{\tabcolsep}{4pt}
\setlength{\LTleft}{0pt}
\setlength{\LTright}{0pt}

\begin{longtable}{
  >{\RaggedRight\hspace{0pt}}p{7.8cm}   % Métrica
  >{\Centering\hspace{0pt}}p{2.2cm}     % Antes
  >{\Centering\hspace{0pt}}p{2.6cm}     % Sprint 1
  >{\Centering\hspace{0pt}}p{2.2cm}     % Cambio
  }

  % --- ENCABEZADO PRIMERA PÁGINA ---
  \caption{Resultados del Sprint 1}
  \label{tab:resultados-sprint1}                                          \\
  \toprule
  \textbf{Métrica} & \textbf{Antes} & \textbf{Sprint 1} & \textbf{Cambio} \\
  \midrule
  \endfirsthead

  % --- ENCABEZADO SIGUIENTES PÁGINAS ---
  \multicolumn{4}{c}{{\bfseries \tablename\ \thetable{} -- Continuación}} \\
  \toprule
  \textbf{Métrica} & \textbf{Antes} & \textbf{Sprint 1} & \textbf{Cambio} \\
  \midrule
  \endhead

  % --- PIE INTERMEDIO ---
  \midrule
  \multicolumn{4}{r}{\textit{Continúa en la siguiente página...}}         \\
  \endfoot

  % --- PIE FINAL ---
  \bottomrule
  \endlastfoot

  Health Score
                   & 5.2
                   & 5.8
                   & +0.6                                                 \\

  velocidad (SP completados)
                   & 38
                   & 42
                   & +10\%                                                \\

  Smells activos
                   & 3
                   & 3
                   & =                                                    \\

  Smells en observación
                   & 0
                   & 1
                   & +1                                                   \\
\end{longtable}
\endgroup

\section{Sprint 2: Abordar smells activos}

\subsection{Contexto}

El conflicto Carlos-Miguel estaba en observación. El foco ahora era el \textit{smell} ``Lone Wolf'' de Ana y la falta de documentación.

\textbf{Social Sprint Backlog}\mbox{}\\[-0.3\baselineskip]

\textbf{Capacidad asignada}: (10 SSP de 50 SP totales = 20\%):
\begin{itemize}
  \item Crear ``Acuerdo de Feedback'' (3 SSP)
  \item Sesiones de transferencia de conocimiento Ana $\rightarrow$ equipo (5 SSP)
  \item Documentar módulo de inventarios (2 SSP)
\end{itemize}

\subsection{Intervención: Transferencia de conocimiento}

Ana facilitó 3 sesiones de 90 minutos cada una donde se habló sobre arquitectura, patrones de integración y debugging de casos críticos. En la Tabla \ref{tab:transferencia-conocimiento} se presentan las sesiones realizadas.

\begingroup
\scriptsize
\renewcommand{\arraystretch}{1.4}
\setlength{\tabcolsep}{4pt}
\setlength{\LTleft}{0pt}
\setlength{\LTright}{0pt}

\begin{longtable}{
  >{\Centering\hspace{0pt}}p{1.8cm}     % Sesión
  >{\RaggedRight\hspace{0pt}}p{9.6cm}   % Tema
  >{\RaggedRight\hspace{0pt}}p{3.8cm}   % Participantes
  }

  % --- ENCABEZADO PRIMERA PÁGINA ---
  \caption{Sesiones de transferencia de conocimiento}
  \label{tab:transferencia-conocimiento}                                                                                                             \\
  \toprule
  \textbf{Sesión} & \textbf{Tema}                                                                                           & \textbf{Participantes} \\
  \midrule
  \endfirsthead

  % --- ENCABEZADO SIGUIENTES PÁGINAS ---
  \multicolumn{3}{c}{{\bfseries \tablename\ \thetable{} -- Continuación}}                                                                            \\
  \toprule
  \textbf{Sesión} & \textbf{Tema}                                                                                           & \textbf{Participantes} \\
  \midrule
  \endhead

  % --- PIE INTERMEDIO ---
  \midrule
  \multicolumn{3}{r}{\textit{Continúa en la siguiente página...}}                                                                                    \\
  \endfoot

  % --- PIE FINAL ---
  \bottomrule
  \endlastfoot

  1
                  & Arquitectura del módulo de inventarios
                  & Todo el equipo                                                                                                                   \\

  2
                  & Patrones de integración con ERP (Enterprise Resource Planning, Planificación de Recursos Empresariales)
                  & Carlos, Pedro                                                                                                                    \\

  3
                  & Debugging de casos críticos
                  & Miguel, Sara                                                                                                                     \\
\end{longtable}
\endgroup

Resultado:
\begin{itemize}
  \item Bus factor del módulo crítico pasó de 1 (solo Ana) a 3 (Ana + Carlos + Miguel).
  \item Ana reportó ``alivio'' por no ser la única responsable.
  \item Se creó wiki con 12 páginas de documentación técnica.
\end{itemize}

\subsection{Evento inesperado: Crisis de deadline}

Día 7: El PO anunció que un cliente importante necesitaba una feature para el día 10 (3 días antes del fin de Sprint).

\begin{escenario}
  \textbf{Product Owner}:
  \begin{quote}
    ``Necesitamos priorizar esta iniciativa. Es crítica para el cliente.''
  \end{quote}
  \textbf{Social Guide (Laura)}:
  \begin{quote}
    ``El Health Score actual es 6.1. El equipo puede absorber presión temporal. Sin embargo, propongo reducir el Social Sprint Backlog a 5 SSP cancelando una sesión de transferencia de conocimiento, para liberar capacidad técnica.''
  \end{quote}
  \textbf{Equipo}:
  \begin{quote}
    Acepta la propuesta.
  \end{quote}
  La sesión de transferencia número 3 se re-programa para el Sprint 3.

\end{escenario}

Este escenario ilustra cómo SocialScrum permite flexibilidad: la salud social no es ``sagrada''; puede ajustarse cuando hay presión legítima, siempre que el Health Score lo permita.

\subsection{Resultados Sprint 2}

En la Tabla \ref{tab:resultados-sprint2} se presentan los resultados cuantitativos del Sprint 2.

\begingroup
\scriptsize
\renewcommand{\arraystretch}{1.4}
\setlength{\tabcolsep}{4pt}
\setlength{\LTleft}{0pt}
\setlength{\LTright}{0pt}

\begin{longtable}{
  >{\RaggedRight\hspace{0pt}}p{7.8cm}   % Métrica
  >{\Centering\hspace{0pt}}p{2.2cm}     % Sprint 1
  >{\Centering\hspace{0pt}}p{2.6cm}     % Sprint 2
  >{\Centering\hspace{0pt}}p{2.4cm}     % Cambio
  }

  % --- ENCABEZADO PRIMERA PÁGINA ---
  \caption{Resultados Sprint 2}
  \label{tab:resultados-sprint2}                                                      \\
  \toprule
  \textbf{Métrica} & \textbf{Sprint 1}          & \textbf{Sprint 2} & \textbf{Cambio} \\
  \midrule
  \endfirsthead

  % --- ENCABEZADO SIGUIENTES PÁGINAS ---
  \multicolumn{4}{c}{{\bfseries \tablename\ \thetable{} -- Continuación}}             \\
  \toprule
  \textbf{Métrica} & \textbf{Sprint 1}          & \textbf{Sprint 2} & \textbf{Cambio} \\
  \midrule
  \endhead

  % --- PIE INTERMEDIO ---
  \midrule
  \multicolumn{4}{r}{\textit{Continúa en la siguiente página...}}                     \\
  \endfoot

  % --- PIE FINAL ---
  \bottomrule
  \endlastfoot

  Health Score
                   & 5.8
                   & 6.4
                   & +0.6                                                             \\

  velocidad
                   & 42
                   & 45
                   & +7\%                                                             \\

  Smells activos
                   & 3
                   & 2
                   & -1                                                               \\

  Smells mitigados
                   & 0
                   & 1 (\textit{Radio Silence})
                   & +1                                                               \\
\end{longtable}
\endgroup

\section{Sprint 3: Consolidación y ajustes}

\subsection{Contexto}

Health Score en 6.4 (zona ``En riesgo'' pero cerca de ``Saludable''). Dos smells activos: \textit{Organisational Silo} (FE-BE) y \textit{Lone Wolf} (reducido a Media).

\textbf{Social Sprint Backlog}\mbox{}\\[-0.3\baselineskip]

\textbf{Capacidad asignada}: (6 SSP de 50 SP = 12\%):
\begin{itemize}
  \item Completar sesión 3 de transferencia (2 SSP).
  \item Ritual de sincronización FE-BE (Frontend-Backend) semanal (2 SSP).
  \item Retrospectiva profunda de 3 Sprints (2 SSP).
\end{itemize}

\subsection{Intervención: Sincronización FE-BE}

Se implementó reunión semanal de 30 minutos entre Frontend y Backend:

\begin{itemize}
  \item Agenda fija: ¿Qué cambios de API vienen? ¿Qué necesita FE de BE? ¿Dónde hay fricción?
  \item Rotación de facilitador (FE una semana, BE la siguiente).
  \item Output: Lista de dependencias para el Sprint Planning.
\end{itemize}

Resultado: Integraciones tardías se redujeron de 4 por Sprint a 1.

\subsection{Retrospectiva de consolidación}

Laura facilitó una retrospectiva extendida (2 horas) para evaluar los resultados obtenidos a lo largo de tres Sprints consecutivos.

\paragraph{\normalfont{Aspectos que funcionaron.}}
\begin{itemize}
  \item El Social Daily Standup permitió detectar tensiones antes de que escalaran.
  \item La sesión de alineación entre Carlos y Miguel evitó el bloqueo de una feature crítica.
  \item La transferencia de conocimiento redujo la dependencia técnica de Ana.
  \item El equipo comenzó a mencionar problemas de forma explícita en retrospectivas.
\end{itemize}

\paragraph{\normalfont{Aspectos que no funcionaron.}}
\begin{itemize}
  \item La sobrecarga inicial se percibió como elevada (la percepción mejoró en los Sprints 2 y 3).
  \item Algunas preguntas del Social Daily Standup se volvieron rutinarias con el tiempo.
  \item El Product Owner mostró resistencia inicial a asignar capacidad a ítems sociales.
\end{itemize}

\paragraph{\normalfont{Ajustes acordados.}}
\begin{itemize}
  \item Reducir el Social Sprint Backlog al 10\% como nivel de mantenimiento.
  \item Rotar la pregunta social semanalmente, en lugar de diariamente.
  \item Incluir al Product Owner en una revisión mensual del Health Score.
\end{itemize}

\subsection{Resultados Sprint 3}

En la Tabla \ref{tab:resultados-sprint3} se presentan los resultados cuantitativos del Sprint 3.

\begingroup
\scriptsize
\renewcommand{\arraystretch}{1.4}
\setlength{\tabcolsep}{4pt}
\setlength{\LTleft}{0pt}
\setlength{\LTright}{0pt}

\begin{longtable}{
  >{\RaggedRight\hspace{0pt}}p{7.8cm}   % Métrica
  >{\Centering\hspace{0pt}}p{2.2cm}     % Sprint 2
  >{\Centering\hspace{0pt}}p{2.6cm}     % Sprint 3
  >{\Centering\hspace{0pt}}p{2.4cm}     % Cambio
  }

  % --- ENCABEZADO PRIMERA PÁGINA ---
  \caption{Resultados Sprint 3}
  \label{tab:resultados-sprint3}                                             \\
  \toprule
  \textbf{Métrica} & \textbf{Sprint 2} & \textbf{Sprint 3} & \textbf{Cambio} \\
  \midrule
  \endfirsthead

  % --- ENCABEZADO SIGUIENTES PÁGINAS ---
  \multicolumn{4}{c}{{\bfseries \tablename\ \thetable{} -- Continuación}}    \\
  \toprule
  \textbf{Métrica} & \textbf{Sprint 2} & \textbf{Sprint 3} & \textbf{Cambio} \\
  \midrule
  \endhead

  % --- PIE INTERMEDIO ---
  \midrule
  \multicolumn{4}{r}{\textit{Continúa en la siguiente página...}}            \\
  \endfoot

  % --- PIE FINAL ---
  \bottomrule
  \endlastfoot

  Health Score
                   & 6.4
                   & 7.1
                   & +0.7                                                    \\

  velocidad
                   & 45
                   & 48
                   & +6\%                                                    \\

  Smells activos
                   & 2
                   & 1
                   & -1                                                      \\

  Smells mitigados
                   & 1
                   & 2
                   & +1                                                      \\
\end{longtable}
\endgroup

\section{Resumen de resultados}

Los resultados presentados son un resumen de la evolución del equipo Phoenix durante los tres Sprints de implementación de SocialScrum. En la Tabla \ref{tab:evolucion-phoenix} se sintetizan las métricas clave, y en la Figura \ref{fig:evolucion-health-score} se visualiza la mejora del Health Score a lo largo del tiempo.

\begingroup
\scriptsize
\renewcommand{\arraystretch}{1.4}
\setlength{\tabcolsep}{4pt}
\setlength{\LTleft}{0pt}
\setlength{\LTright}{0pt}

\begin{longtable}{
  >{\RaggedRight\hspace{0pt}}p{6.0cm}   % Métrica
  >{\Centering\hspace{0pt}}p{2.0cm}     % Inicial
  >{\Centering\hspace{0pt}}p{2.2cm}     % Sprint 1
  >{\Centering\hspace{0pt}}p{2.2cm}     % Sprint 2
  >{\Centering\hspace{0pt}}p{2.2cm}     % Sprint 3
  }

  % --- ENCABEZADO PRIMERA PÁGINA ---
  \caption{Evolución del equipo Phoenix en tres Sprints}
  \label{tab:evolucion-phoenix}                                                                   \\
  \toprule
  \textbf{Métrica} & \textbf{Inicial} & \textbf{Sprint 1} & \textbf{Sprint 2} & \textbf{Sprint 3} \\
  \midrule
  \endfirsthead

  % --- ENCABEZADO SIGUIENTES PÁGINAS ---
  \multicolumn{5}{c}{{\bfseries \tablename\ \thetable{} -- Continuación}}                         \\
  \toprule
  \textbf{Métrica} & \textbf{Inicial} & \textbf{Sprint 1} & \textbf{Sprint 2} & \textbf{Sprint 3} \\
  \midrule
  \endhead

  % --- PIE INTERMEDIO ---
  \midrule
  \multicolumn{5}{r}{\textit{Continúa en la siguiente página...}}                                 \\
  \endfoot

  % --- PIE FINAL ---
  \bottomrule
  \endlastfoot

  Health Score
                   & 5.2
                   & 5.8
                   & 6.4
                   & 7.1                                                                          \\

  velocidad (SP)
                   & 38
                   & 42
                   & 45
                   & 48                                                                           \\

  Smells activos
                   & 3
                   & 3
                   & 2
                   & 1                                                                            \\

  Capacidad social (\%)
                   & 0\%
                   & 16\%
                   & 20\%
                   & 12\%                                                                         \\
\end{longtable}
\endgroup

\begin{figure}[h]
  \centering
  \begin{tikzpicture}[scale=0.9]
    % Ejes
    \draw[thick, -stealth] (0,0) -- (10,0) node[right, font=\footnotesize] {Sprints};
    \draw[thick, -stealth] (0,0) -- (0,6) node[above, font=\footnotesize] {Health Score};

    % Lineas de zona
    \fill[red!10] (0,0) rectangle (10,2.45);
    \fill[yellow!10] (0,2.45) rectangle (10,3.45);
    \fill[green!10] (0,3.45) rectangle (10,5);

    % Etiquetas de zona
    \node[font=\tiny, gray] at (0.8, 1.2) {Critico};
    \node[font=\tiny, gray] at (0.8, 2.9) {Riesgo};
    \node[font=\tiny, gray] at (0.8, 4.2) {Saludable};

    % Escala Y
    \node[font=\tiny] at (-0.3, 0) {1};
    \node[font=\tiny] at (-0.3, 2.45) {5};
    \node[font=\tiny] at (-0.3, 3.45) {7};
    \node[font=\tiny] at (-0.3, 5) {10};

    % Puntos de datos (escalados: 5.2->2.6, 5.8->2.9, 6.4->3.2, 7.1->3.55)
    \fill[blue!70] (1,2.6) circle (0.12) node[above, font=\tiny] {5.2};
    \fill[blue!70] (3.5,2.9) circle (0.12) node[above, font=\tiny] {5.8};
    \fill[blue!70] (6,3.2) circle (0.12) node[above, font=\tiny] {6.4};
    \fill[blue!70] (8.5,3.55) circle (0.12) node[above, font=\tiny] {7.1};

    % Linea de conexión
    \draw[thick, blue!70] (1,2.6) -- (3.5,2.9) -- (6,3.2) -- (8.5,3.55);

    % Etiquetas de Sprint
    \node[font=\tiny] at (1, -0.4) {Inicial};
    \node[font=\tiny] at (3.5, -0.4) {Sprint 1};
    \node[font=\tiny] at (6, -0.4) {Sprint 2};
    \node[font=\tiny] at (8.5, -0.4) {Sprint 3};

    % Flecha de tendencia
    \draw[-stealth, thick, green!60!black] (9, 2.6) -- (9, 3.55);
    \node[font=\tiny, green!60!black, align=center] at (9.5, 3.1) {+1.9\\puntos};
  \end{tikzpicture}
  \caption{Evolución del Health Score del equipo Phoenix}
  \label{fig:evolucion-health-score}
\end{figure}

\section{Lecciones aprendidas}

Durante la implementación de SocialScrum en el equipo Phoenix, se identificaron varios factores clave que contribuyeron al éxito del enfoque, así como dificultades que surgieron y recomendaciones para otros equipos interesados en adoptar prácticas similares.

\subsection{Factores de éxito}

\begin{enumerate}
  \item \textbf{Compromiso visible del liderazgo}: El CTO aprobó dedicar hasta 20\% de capacidad a items sociales en los primeros Sprints.

  \item \textbf{Intervención temprana en conflicto}: La sesión Carlos-Miguel en Sprint 1 evitó que el conflicto bloqueara una feature crítica.

  \item \textbf{Flexibilidad ante presión}: Sprint 2 demostró que SocialScrum puede adaptarse a deadlines urgentes sin abandonar completamente la agenda social.

  \item \textbf{Transferencia de conocimiento proactiva}: Reducir el bus factor de Ana benefició tanto la salud social como la resiliencia técnica.
\end{enumerate}

\subsection{Dificultades encontradas}

\begin{enumerate}
  \item \textbf{Resistencia inicial del PO}: Percibía los SSP como ``tiempo robado'' al desarrollo. Se resolvió mostrando correlación velocidad-Health Score.

  \item \textbf{Fatiga de preguntas sociales}: El equipo se ``acostumbró'' a las preguntas del Daily. Se resolvió rotando preguntas semanalmente.

  \item \textbf{Dificultad para cuantificar ROI}: Es difícil probar causalidad entre intervenciones sociales y mejora de velocidad. Se aceptó correlación como evidencia suficiente.
\end{enumerate}

\subsection{Recomendaciones para otros equipos}

En la Tabla \ref{tab:recomendaciones-phoenix} se resumen recomendaciones prácticas basadas en la experiencia del equipo Phoenix.

\begingroup
\scriptsize
\renewcommand{\arraystretch}{1.4}
\setlength{\tabcolsep}{4pt}
\setlength{\LTleft}{0pt}
\setlength{\LTright}{0pt}

\begin{longtable}{
  >{\RaggedRight\hspace{0pt}}p{5.5cm}   % Recomendación
  >{\RaggedRight\hspace{0pt}}p{10.0cm}   % Justificación
  }

  % --- ENCABEZADO PRIMERA PÁGINA ---
  \caption{Recomendaciones basadas en la simulación del equipo Phoenix}
  \label{tab:recomendaciones-phoenix}                                                                                                                \\
  \toprule
  \textbf{Recomendación} & \textbf{Justificación}                                                                                                    \\
  \midrule
  \endfirsthead

  % --- ENCABEZADO SIGUIENTES PÁGINAS ---
  \multicolumn{2}{c}{{\bfseries \tablename\ \thetable{} -- Continuación}}                                                                            \\
  \toprule
  \textbf{Recomendación} & \textbf{Justificación}                                                                                                    \\
  \midrule
  \endhead

  % --- PIE INTERMEDIO ---
  \midrule
  \multicolumn{2}{r}{\textit{Continúa en la siguiente página...}}                                                                                    \\
  \endfoot

  % --- PIE FINAL ---
  \bottomrule
  \endlastfoot

  Empezar con 15--20\% de capacidad social
                         & Permite intervenciones significativas sin colapsar la velocidad. Reducir a 10\% una vez estabilizado el equipo.           \\

  Priorizar conflictos activos
                         & Son los factores que más bloquean el trabajo. Resolverlos temprano libera capacidad para acciones preventivas.            \\

  Medir Health Score semanalmente
                         & Permite detectar deterioro antes de que se convierta en crisis. La medición mensual es insuficiente en equipos en riesgo. \\

  Involucrar al PO desde el Sprint 0
                         & La resistencia inicial es predecible. La educación temprana facilita entender la correlación entre salud y productividad. \\

  Documentar intervenciones
                         & Permite aprender qué prácticas funcionan para el equipo específico y justificar inversiones futuras.                      \\
\end{longtable}
\endgroup

La simulación del equipo Phoenix demuestra que SocialScrum puede implementarse de forma incremental, con resultados medibles en 3 Sprints. La mejora del Health Score de 5.2 a 7.1 (+36\%) coincidió con un aumento de velocidad de 38 a 48 SP (+26\%). Aunque la correlación no prueba causalidad, la consistencia de ambas tendencias sugiere que abordar la deuda social tiene beneficios tanto para el bienestar del equipo como para su productividad.
